\Needspace{18\baselineskip}
\fancyhead[L]{\small COMET OF 1807.}
\begin{center}
COMET OF 1807.
\end{center}

\begin{center}
\fontsize{10}{12}\selectfont\begin{tabular}{c}
\hline
Göttingische gelehrte Anzeigen. Issue 32. Pp. 313..315. 25 February 1808. \\
\hline
\end{tabular}
\end{center}

The unfavorable weather throughout the winter allowed only a few observations of the latest comet at our observatory, most of them under poor conditions. Since these observations will by now have ended everywhere, we present here all those that could be made here since the arrival of Prof. \textsc{Gauss}.


\begin{center}\fontsize{9}{11}\selectfont\setlength{\tabcolsep}{2pt}\renewcommand{\arraystretch}{1.15}\begin{tabular}{lrrrrr}
& & \multicolumn{2}{c}{Apparent} & \multicolumn{2}{c}{Difference}\\
& Mean time & \shortstack{right\\ascension} & \shortstack{north\\declination} & \shortstack{right\\asc.} & declination\\\hline
1807 Dec. 17. & $11^{\text{u}}\,24^{\text{m}}\,44^{\text{s}}$ & $305^\circ\;54'\;0''$ & $45^\circ\;4'\;30''$ & $+49$ & $-28$\\
29 & $7\;59\;21$ & $318\;57\;41$ & $46\;41\;48$ & $+6$ & $-25$\\
29 & $11\;36\;33$ & $319\;6\;35$ & $46\;43\;11$ & $+49$ & $-55$\\
1808 Jan. 4 & $7\;22\;12$ & $325\;11\;19$ & $47\;12\;46$ & $-18$ & $-21$\\
4 & $9\;13\;57$ & $325\;15\;53$ & & $-6$ & \\
4 & $17\;30\;26$ & $325\;36\;38$ & & $+18$ & \\
31 & $9\;31\;34$ & $349\;46\;43$ & $48\;9\;26$ & $-84$ & $+24$\\
\end{tabular}\end{center}

\par\indent All these observations except the last were made with the ring micrometer by comparison with stars in the \emph{histoire céleste}. For the last observation a seventh-magnitude star whose position had yet to be determined was used for comparison. An observation of its position could only be obtained on 14 February. It gave


\begin{center}\fontsize{10}{12}\selectfont\setlength{\tabcolsep}{2pt}\renewcommand{\arraystretch}{1.15}\begin{tabular}{lr}
apparent right ascension of the star .\;.\;.\;.\; & $350^\circ\;14'\;38''$\\
apparent northern declination .\;.\;.\;.\;.\;.\;.\; & $48\;4\;33$\\
\end{tabular}\end{center}

\noindent The observation above was reduced using this position, which still needs to be corrected.

\par\indent The parabolic elements of the comet, which Prof. \textsc{Gauss} had already calculated at the beginning of November and which are known from the Monatliche Correspondenz, agreed with these later observations as follows [the differences have been added to the observations above in this reprint.]

\par\indent It would scarcely be worth improving those elements again from these observations, especially since the comet was so faint that its entries and exits on 31 January could only be observed with difficulty, and the last observations could therefore hardly be very accurate. Unless observers at other observatories had much better luck, it will scarcely be possible to say anything decisive about how far the comet’s orbit departs from a parabola, despite its long visi\SourceLeaf{304}bility. In any event, the departure can certainly only be very small.

\par\indent On 14 February the comet was still just visible in the \textsc{Herschel} telescope, but a proper observation was out of the question.
